% ============================================================================
%  Actividad 4 · Anexo
%
%  Propiedad de US-UX-07. Reune el protocolo de captura, el inventario de
%  figuras y de laminas, el inventario tecnico de lo construido, la procedencia
%  de los colores del tema institucional y la trazabilidad entre los criterios
%  de la historia de usuario, las secciones de este documento y su evidencia.
%
%  Los codigos de color de la tabla de procedencia se escriben SIN el signo de
%  numero a proposito: scripts/verificar_tokens_a4.sh falla si un archivo de
%  contenido contiene ese patron. La tabla lo explica en su propio texto.
% ============================================================================

% ============================================================================
\uxsection{Anexo}

\subsection{Protocolo de captura, en resumen}

El guion completo, con sus precondiciones y su ruta manual alternativa, está archivado en el
repositorio del proyecto junto al guion ejecutable que lo implementa. Lo que sigue es su resumen
operativo, suficiente para reproducir cualquier figura de este documento.

\begin{uxtabla}{|L{4.2cm}|Y|}{Parámetros del protocolo de captura}
  \uxheadrow \thd{Parámetro} & \thd{Valor} \\
  Resolución de composición & 1440 por 900 píxeles \\
  Modo de color y tema & Modo claro y tema de omisión, fijados explícitamente antes de la primera captura \\
  Idioma de la interfaz & Español \\
  Sesión & Una por pantalla, con el rol que el contrato de navegación sugiere, abierta desde la propia interfaz \\
  Momento del disparo & Con la red inactiva y la franja de alcance presente en la página \\
  Nombre de archivo & Número de pantalla, ruta y estado, separados por guion bajo \\
  Origen del plan & El contrato de navegación del portal, leído por el guion; no una lista escrita aparte \\
\end{uxtabla}

El nombre de archivo es el contrato que empareja el antes y el después: las dos carpetas usan los
mismos nombres y la pareja se detecta por igualdad. Una captura del después con otro nombre no forma
pareja con nada, y la iteración documentada deja de ser verificable.

\subsection{Inventario de figuras}

\begin{uxtablalarga}{|C{0.8cm}|L{5.7cm}|L{2.4cm}|Y|}{Inventario de las figuras de esta actividad}{\thd{\#} & \thd{Archivo} & \thd{Sesión} & \thd{Dónde se usa}}
  0 & \texttt{0\_acceso\_normal.png} & sin sesión & Prototipo 0 \\
  1 & \texttt{1\_inicio\_normal.png} & operativo & Prototipo 1 \\
  2 & \texttt{2\_exploracion\_normal.png} & analista & Prototipo 2, y el par antes y después de la iteración \\
  3 & \texttt{3\_gobierno\_normal.png} & analista & Prototipo 3 \\
  4 & \texttt{4\_asistente\_normal.png} & directivo & Prototipo 4 \\
  5 & \texttt{5\_administracion\_normal.png} & administración & Prototipo 5 \\
  6 & \texttt{6\_exploracion-exportar\_normal.png} & analista & Prototipo 6 \\
  4 & \texttt{4\_asistente\_resultado.png} & directivo & Prototipo 4, turno resuelto con su tarjeta de llamada a herramienta \\
\end{uxtablalarga}

Los siete nombres del estado normal existen dos veces, en la carpeta del antes y en la del después.
Las figuras de la sección de prototipos usan las del después, que corresponden al estado entregado;
la sección de pre-validación usa la pareja completa de la pantalla afectada por la iteración.

La octava figura, el turno resuelto del asistente, existe solo en la carpeta del después y de manera
deliberada: documenta un estado de la conversación y no una versión distinta de la pantalla, de modo
que emparejarla con un antes sugeriría un cambio que no ocurrió.

\subsection{Inventario de láminas del archivo de diseño}

Además de las capturas, la entrega se apoya en las láminas exportadas del archivo de trabajo y
archivadas en el repositorio. Se listan las veintiuna, por la página del archivo a la que
pertenecen, de modo que cualquiera pueda localizar en el archivo vivo lo que ve impreso aquí.

\begin{uxtablalarga}{|L{2.1cm}|L{6.5cm}|Y|}{Láminas exportadas del archivo de diseño}{\thd{Página} & \thd{Archivo} & \thd{Qué muestra}}
  Guía de estilo & \texttt{guia\_logotipo.png} & Construcción del logotipo, resguardo y usos incorrectos \\
  Guía de estilo & \texttt{guia\_tipografia.png} & Espécimen y escala de roles tipográficos \\
  Guía de estilo & \texttt{guia\_paleta.png} & Paleta de la identidad con sus papeles \\
  Guía de estilo & \texttt{guia\_cuadricula.png} & Retícula, espaciado y comportamiento adaptativo \\
  Guía de estilo & \texttt{guia\_iconografia.png} & Familia de iconos, tamaños y grosores \\
  Guía de estilo & \texttt{guia\_versionado.png} & Regla de documentación y control de versiones del archivo \\
  Guía de estilo & \texttt{guia\_interaccion\_accesibilidad\_voz.png} & Microinteracciones, criterios de accesibilidad y voz del producto \\
  Componentes & \texttt{componente\_botones.png} & Jerarquía de botones con sus estados \\
  Componentes & \texttt{componente\_busqueda.png} & Los cinco estados del campo de búsqueda \\
  Componentes & \texttt{componente\_tabla.png} & Tabla de datos, cabecera, densidad y alineación de cifras \\
  Componentes & \texttt{componente\_alertas.png} & Alertas y mensajes por semántica \\
  Prototipos & \texttt{flujo\_1\_variables.png} & Pantalla protagonista del flujo de análisis y exportación \\
  Prototipos & \texttt{flujo\_2\_riesgo.png} & Pantalla protagonista del flujo de señal de riesgo \\
  Prototipos & \texttt{flujo\_3\_versiones.png} & Pantalla protagonista del flujo de publicación de versiones \\
  Prototipos & \texttt{flujo\_4\_autorizacion.png} & Pantalla protagonista del flujo de autorización de accesos \\
  Prototipos & \texttt{flujo\_5\_centro\_trabajo.png} & Pantalla protagonista del flujo del centro de trabajo \\
  Prototipos & \texttt{secuencia\_flujo\_1.png} & Las ocho pantallas del flujo de análisis y exportación, en orden \\
  Prototipos & \texttt{secuencia\_flujo\_2.png} & Las cinco pantallas del flujo de señal de riesgo, en orden \\
  Prototipos & \texttt{secuencia\_flujo\_3.png} & Las seis pantallas del flujo de publicación de versiones, en orden \\
  Prototipos & \texttt{secuencia\_flujo\_4.png} & Las cinco pantallas del flujo de autorización de accesos, en orden \\
  Prototipos & \texttt{secuencia\_flujo\_5.png} & Las seis pantallas del flujo del centro de trabajo, en orden \\
\end{uxtablalarga}

Las láminas que este documento imprime salen de esta lista; las demás quedan como material de
consulta del archivo, disponible por el enlace que publica la sección de método.

\subsection{Inventario técnico de lo construido en esta actividad}

El inventario siguiente declara qué se construyó y con qué se comprueba cada cosa. La columna de la
derecha no dice \enquote{se revisó}: nombra el mecanismo que vuelve a comprobarlo cada vez que
alguien corre la suite.

\begin{uxtablalarga}{|L{4.2cm}|L{4.6cm}|Y|}{Lo construido en esta actividad y su mecanismo de comprobación}{\thd{Pieza} & \thd{Qué es} & \thd{Cómo se comprueba}}
  Eje de tema del sistema de diseño & Los diecisiete tokens de color resueltos para dos temas y dos modos, con la familia tipográfica como parte del eje & Cálculo de contraste sobre las cuatro combinaciones: cero tokens que informan por debajo de su listón \\
  Fijación del tema de omisión & Los diecisiete valores del tema de omisión, en sus dos modos & Prueba escrita antes de abrir el eje, que falla si un solo valor cambia. Es lo que garantiza que las capturas ya entregadas sigan correspondiendo al producto \\
  Separación bajo dicromacia & El peor par de colores semánticos bajo las tres dicromacias, por tema y modo & Piso de 13{,}6 en claro y 21{,}5 en oscuro, con 14{,}5 y 22{,}6 medidos en el tema institucional \\
  Emisión de los tokens & La hoja de estilo del portal y su paleta tipada, generadas desde la definición única & Dos corridas seguidas del generador producen archivos byte a byte iguales, y un guion compara el disco contra lo que el generador produce hoy \\
  Conmutador de tema y conmutador de perfil & Dos controles en la cabecera: el tema como preferencia que persiste, el perfil como sesión que se acuña de verdad contra el portal & Pruebas de montaje sobre los dos controles, incluido el caso de la puerta de demostración apagada, en el que el conmutador de perfil no se dibuja \\
  Entrada en frío desde el índice & La tarjeta de cada prototipo abre su pantalla con el perfil que esa pantalla sugiere, en un solo movimiento; sin sesión, el desvío dice a dónde se iba y devuelve ahí & Pruebas sobre la decisión de la guarda y sobre el destino de retorno, que se valida contra el contrato de navegación \\
  Pantalla de catálogo & Buscador, conteos por dominio, resultados y los cuatro estados no felices, compuestos desde la búsqueda del catálogo & Once casos de prueba que montan la pantalla y sus estados, y la comprobación de que el componente de andamiaje ya no aparece en ella \\
  Etiquetas de alcance & Seis corregidas y una conservada, según la regla publicada en la sección de alcance & Prueba que lee la tabla de este documento y la compara, fila por fila, contra el contrato de navegación del portal \\
\end{uxtablalarga}

\subsection{Procedencia de los colores del tema institucional}

El sistema de diseño del portal resuelve dos temas. El de omisión es el que sostiene las capturas de
este documento; el institucional toma su color del archivo de diseño. Esta tabla declara, color por
color, qué valor llegó del archivo y cuál es derivación, con la regla que la produjo. Sin esta
distinción, una paleta se lee como si todos sus valores tuvieran el mismo respaldo, y no lo tienen.

\begin{uxtablalarga}{|L{2.2cm}|L{1.6cm}|L{3.2cm}|Y|}{Procedencia de los colores del tema institucional}{\thd{Color del archivo} & \thd{Código} & \thd{Token que lo recibe} & \thd{Regla aplicada}}
  Superficie & \texttt{FFFFFF} & Suelo del modo claro & Se adopta sin cambio \\
  Navegación & \texttt{102A43} & Corriente plena en claro, y superficie de panel en oscuro & Se adopta sin cambio. El suelo del modo oscuro, \texttt{0B1B2B}, se deriva de este color oscureciéndolo, y no del negro puro: un tema institucional no se apaga en gris \\
  Secundario & \texttt{1D4C6E} & Corriente media en claro & Se adopta sin cambio \\
  Éxito & \texttt{287A58} & Confirmación & Sin cambio en el modo claro; aclarado para el oscuro \\
  Atención & \texttt{B97812} & Aviso & Oscurecido un 12\,\% conservando matiz y saturación, hasta \texttt{A36A10}. El valor del archivo daba 3{,}65:1 y no alcanzaba el listón de un token que informa \\
  Error & \texttt{B8443F} & Error & Oscurecido un 20\,\% hasta \texttt{933632} en el modo claro; aclarado para el oscuro y oscurecido después un 12\,\% hasta \texttt{EC5B51}. Los valores intermedios se separaban poco de la confirmación bajo protanopia \\
  Ninguno & --- & Canal informativo & Ningún matiz del octeto de la marca se separa de la confirmación bajo tritanopia, así que el valor sale del octeto: azul profundo en el modo claro y violeta en el oscuro, donde el azul y el verde colapsan \\
  Ninguno & --- & Rampa de neutros y superficies auxiliares & Derivados del color de navegación por variación de luminosidad, con la razón de contraste de cada paso verificada por cálculo \\
  Ninguno & --- & Las seis series de datos & Conservan el valor del tema de omisión: son canal de datos y no identidad, y su separación ya estaba verificada \\
\end{uxtablalarga}

Los códigos de esta tabla se imprimen sin el signo de número que suele precederlos, y no es un
descuido de composición. Una comprobación del repositorio recorre los archivos de contenido de esta
actividad buscando exactamente ese patrón, para impedir que un color se teclee a mano en el
documento en lugar de salir de su generador. Los valores de arriba describen la procedencia de una
paleta que vive en el sistema de tokens del portal ---otra cadena, con otro generador--- y por eso
se citan en esta forma, que la comprobación distingue de un color escrito a mano.

\subsection{Trazabilidad de los criterios de aceptación}

\begin{uxtablalarga}{|L{1.6cm}|L{5.2cm}|Y|}{Dónde se atiende cada criterio y con qué evidencia}{\thd{Criterio} & \thd{Enunciado} & \thd{Sección y evidencia}}
  1 & Siete prototipos navegables y desplegados & Sección de prototipos, siete figuras con pie numerado \\
  2 & Sistema de diseño de fuente única, con derivación en un solo sentido & Sección de método, cadena de derivación; guía de estilos, láminas emitidas por el generador \\
  3 & Cada ruta anclada a una rama del mapa & Sección de prototipos, tabla de veinte filas verificada por prueba automática \\
  4 & Los cuatro estados no felices por pantalla & Sección de prototipos, matriz de veintiocho celdas y las nueve reglas de comportamiento de las celdas especificadas \\
  5 & Tabla de alcance por pantalla y por historia & Sección de alcance, con la regla de asignación escrita, la tabla de las siete pantallas y el catálogo completo de historias. El esquema desdobla la hoja de ruta en diseñada y sin diseñar \\
  6 & Pre-validación con iteración documentada & Sección de pre-validación: tres iteraciones con su hallazgo, su cambio y su versión, y el par de figuras de la primera \\
  7 & Nota de herramienta en la introducción & Introducción, cuatro propiedades de la rúbrica y objetivo de aprendizaje \\
  8 & Entrega en tiempo con portada de los tres integrantes & Portada de este documento \\
  9 & Acceso al archivo de trabajo & Sección de método, con el enlace, la estructura de las tres páginas y sus nodos, y lo que se recorre desde cada una \\
  10 & Segunda paleta verificada con el mismo listón & Sección de pre-validación, registro de la segunda iteración; paleta completa en la sección del tema institucional; procedencia de cada valor en este anexo \\
\end{uxtablalarga}

Los dos últimos criterios se incorporaron al consolidar la entrega y se anotan aquí con el mismo
formato que los ocho anteriores, para que la tabla siga sirviendo de guía de lectura completa.

\subsection{Verificaciones automáticas asociadas a este documento}

Dos afirmaciones de este documento no dependen de que alguien las revise a mano: una prueba
automática las compara contra el código del portal en cada corrida de la suite.

\begin{uxtabla}{|L{5.4cm}|Y|}{Lo que una prueba automática verifica de este documento}
  \uxheadrow \thd{Afirmación} & \thd{Defecto que la prueba impide} \\
  La tabla de correspondencia publica las dieciséis ramas del mapa y las cuatro categorías & Que la tabla, escrita a mano, omita una rama y el documento afirme una cobertura que el contrato de navegación no tiene \\
  El alcance publicado por pantalla coincide con el que el código declara & Que el código y el documento diverjan tras el congelamiento, y el entregable calificado describa un producto distinto del entregado \\
\end{uxtabla}

Las dos pruebas leen este documento como entrada. Es una dependencia deliberada: la alternativa
---comparar el documento consigo mismo--- pasa siempre, diga lo que diga.

\subsection{Límites de esta entrega}

Se enuncian aquí para que no haya que deducirlos.

\begin{uxlist}
  \lead{Los evaluadores de la pre-validación son sintéticos}{Sus hallazgos orientan decisiones de
    presentación y no cierran ninguna pregunta sobre las personas usuarias. La primera evidencia
    humana llega con la prueba de usabilidad de la Actividad 5.}
  \lead{La estabilidad visual durante la carga se declara, no se mide}{El entorno del proyecto no
    incluye una herramienta que produzca esa métrica, y una cifra sin instrumento sería peor que
    declarar el método de verificación.}
  \lead{El contenido del asistente está guionizado}{El transporte, las tarjetas de llamada a
    herramienta y la cancelación son reales. La conexión con un modelo de lenguaje queda declarada
    en hoja de ruta y la pantalla lo dice en su propia franja.}
  \lead{El catálogo construido no tiene captura en este documento}{Su estado entregado se
    describe en la sección de prototipos y se comprueba con pruebas automáticas. La única captura de
    esa ruta sostiene el par de antes y después de la primera iteración, que es evidencia que no se
    puede volver a tomar.}
  \lead{Las láminas del archivo de diseño no son capturas del portal}{Documentan decisiones de
    interfaz tomadas, no capacidades construidas. La sección de alcance las usa exactamente para
    eso: distinguir lo que está diseñado y sin construir de lo que ni siquiera está diseñado.}
\end{uxlist}
